\begin{abstract}
This is the technical companion to the conceptual account of Vibe Theory. Where the conceptual
paper tells the story in plain terms, this one walks through the machinery and the findings: the
finite, seeded testbed, the geometry of the substrate, and the eighty-three known-answer
experiments that build the model from the bottom up. The substrate is the centerpiece. It is a
growing hyperbolic crystal, a regular honeycomb of a Coxeter reflection group, and we treat it in
detail: the Coxeter construction and the regular tessellations $\{p,q\}$ of the hyperbolic plane,
the heptagrid $\{7,3\}$, the regular honeycombs of hyperbolic space and the dodecagrid $\{5,3,4\}$
that is the model's real three-dimensional substrate, the classification that forces the spatial
dimension into the window two-to-four and then to three, the Lorentz safety that comes from
curvature rather than disorder, and the deterministic automaton that grows the crystal one cell at
a time, forever. On that substrate sit the ternary tones and one local rule, from which we recover
emergent geometry and time, the field and gauge content, a conserving Einstein equation and a
discrete graviton, the quantum pillars with a derived Born rule, a cosmology with a dark sector,
and sharp predictions that pass current bounds while excluding a lattice. We add the Standard Model
results, anomaly-forced charge quantization, the Sakharov conditions, and the mass hierarchy from
hyperbolic overlaps, and we mark honestly what is open. This is a simulable model with a concrete
build target, not a finished empirical theory.
\medskip

\noindent\textbf{Status and caveats.} This is exploratory work, a feasibility study rather than a validated result. The aim is to see whether a model of this kind can be made to hang together at all, and so far it looks promising. Much validation and further work remain. Several of the experiments reported here are preliminary, and some may turn out to be inaccurate or improperly set up. Nothing here should be read as established or taken too rigorously at this stage. It is offered in the spirit of curiosity and invitation, a direction that looks worth exploring rather than a finished theory or a claim of discovery.
\end{abstract}
